\documentclass{tietreport}

\coursecode{UCS503 - Software Engineering}
\projecttitle{SV Car Customs}
\projectsubtitle{An E-Commerce Platform for Vehicle-Compatible Car Accessories}
\students{%
  \texttt{1024170242} Vishu Verma\\
  \texttt{1024170241} Raghu Verma\\
  \texttt{1024170229} Pratham Sharma\\
  \texttt{1024170231} Tarush Mohindru}
\group{Group: \texttt{3Q24} BE Third Year, COPC}
\supervisor{Ms.\ Anushka}
\evaluation{Final Evaluation}

\begin{document}
\maketitle
\pagenumbering{roman}
\tableofcontents
\listoffigures
\listoftables
\cleardoublepage
\pagenumbering{arabic}

% =====================================================================
\chapter{Introduction}
% =====================================================================

\section{Background and Context}

SV Car Customs is a Delhi-based business that sells aftermarket car
accessories (body kits, spoilers, diffusers, exhausts, lighting, steering
wheels, paddle shifters and interior parts) to customers across India. The
business previously sold through a hosted no-code storefront, which gave
little control over search visibility, checkout, payment handling and product
data.

This project replaces that storefront with a custom platform that the team
designed, built, tested and deployed. The platform is live at
\url{https://www.svcarcustoms.in} and is used by the business for real
orders. It has two parts: a customer storefront, where shoppers browse by car
model and part type, check vehicle fitment and pay online, and an
administration panel, where the owner manages the catalogue, orders, settings
and marketing content without developer help.

This final report documents the complete system: its requirements, its design
(UML, DFD and data model), its implementation, how it was verified, and what
was measured after it went live. All figures in Chapter~\ref{ch:results} were
measured on the code submitted with this report (repository revision of
26~November~2026) or on the live site.

\subsection{Problem Statement}

A generic product catalogue cannot express whether an accessory fits a
particular car make and model, so customers order parts that do not fit or
abandon the purchase. Customers also need correct prices and a payment flow
that cannot be manipulated from the browser, while the owner needs one place
to maintain products, variants, stock, orders and Cash-on-Delivery (COD)
collection.

The project therefore addresses the need for a \emph{vehicle-aware commerce
system} in which:
\begin{itemize}
  \item products are linked to the vehicles they fit, and shoppers can filter
        the catalogue by their own car;
  \item the server, not the browser, computes every price and confirms every
        payment;
  \item no paid order is lost even if the customer closes the browser during
        payment; and
  \item the owner controls catalogue, prices, settings and fulfilment from an
        admin panel.
\end{itemize}

\section{Project Objectives}

The objectives below were fixed in the Software Requirements Specification
(SRS, version 1.0, 15 October 2026) and are written in SMART form.
Chapter~\ref{ch:results} reports which ones were met.

\begin{enumerate}
  \item \textbf{Objective 1 -- Fitment-aware shopping.} A customer can find a
        part that fits their car and finish a purchase on a phone in at most
        three checkout screens, by the end of the semester.
  \item \textbf{Objective 2 -- Owner control.} The owner can change products,
        prices, categories, home-page content, tax and shipping settings
        without a code change, and every admin change is recorded in an
        audit log.
  \item \textbf{Objective 3 -- Payment integrity.} Every online payment is
        confirmed by a server-side hash check, and a background job settles
        orders whose payment callback never arrived, so no paid order is lost.
  \item \textbf{Objective 4 -- Search visibility.} The store publishes a
        sitemap, product structured data and a Google Merchant Center feed,
        and every public page is indexable.
  \item \textbf{Objective 5 -- Quality.} Automated tests run on every change
        through continuous integration (CI); pages are served from cache with
        a target Largest Contentful Paint (LCP) under 2.5\,s on 4G.
\end{enumerate}

% =====================================================================
\chapter{Methodology and Design}
% =====================================================================

\section{Software Development Life Cycle}

The team followed an \textbf{iterative and incremental} model with short
(roughly weekly) iterations, which matches the course requirement of
``build, test and ship'' in fast cycles. A pure Waterfall model was rejected
because requirements such as payment-gateway behaviour, image hosting limits
and search ranking could only be learnt by shipping and observing. Each
increment followed the same loop: plan $\rightarrow$ design $\rightarrow$
implement $\rightarrow$ test in CI $\rightarrow$ deploy $\rightarrow$ review.

\begin{table}[H]
\centering\small
\begin{tabularx}{\linewidth}{l X}
\toprule
\textbf{Increment} & \textbf{Delivered scope} \\
\midrule
Phase 0 -- Foundations & Monorepo, database schema, migrations, contracts package, deployment. \\
Phase 1 -- Admin & Admin login, product/category management, image upload. \\
Phase 2 -- Read paths & Catalogue, category and product pages, server-side rendering. \\
Phase 3 -- Auth and cart & Email OTP login, JWT sessions, persistent cart. \\
Phase 4 -- Payments & Checkout, payment gateway hand-off, webhooks (moved to PayU). \\
Phase 5 -- Caching and observability & Redis cache, rate limiting, Sentry, structured logs. \\
Phase 6 -- Search and polish & Search, vehicle fitment, ``My Garage'', SEO. \\
Post-launch increments & Guest checkout, COD advance, promo codes, analytics, CI guards, mobile fixes. \\
\bottomrule
\end{tabularx}
\caption{Development increments (from the project's phase documents and git history)}
\label{tab:phases}
\end{table}

The work spans 118 commits on the main branch between 31~May and
25~November~2026. Changes reach production only through pull requests that
pass CI (Section~\ref{sec:testing}).

\section{Requirements Summary}

The full SRS (version 1.1) is submitted separately (\texttt{srs/main.pdf} in
the repository). It defines 47 functional requirements in seven groups
(catalogue, search and fitment, cart and checkout, payments, accounts,
admin, SEO) and 15 non-functional requirements (performance, reliability,
security). Table~\ref{tab:trace} traces the high-priority requirements to
their implementation and to the evidence used to verify them.

\begin{table}[H]
\centering\footnotesize
\begin{tabularx}{\linewidth}{l X X}
\toprule
\textbf{Requirement} & \textbf{Implemented in} & \textbf{Verified by} \\
\midrule
FR-CAT-02/06 Category pages, hide drafts & \texttt{/category/[slug]}, products API filters on \texttt{status=active} & Storefront and API unit tests; live page \\
FR-SRC-02/03 Fitment & \texttt{vehicle\_makes}, \texttt{vehicle\_models}, \texttt{product\_fitments}; PDP fitment list & API tests; live PDP \\
FR-CHK-04 Server pricing & \texttt{lib/pricing.ts}, \texttt{lib/place-order.ts} & Pricing unit tests \\
FR-CHK-05 Guest checkout & \texttt{/guest-checkout} route & API tests \\
FR-PAY-01/02 Signed hash, verify & \texttt{lib/payu.ts}, \texttt{/webhooks} & Hash and webhook tests \\
FR-PAY-03 COD advance & \texttt{lib/payment-amount.ts} (15\% advance) & Payment-amount tests \\
FR-PAY-05 Reconciliation & Inngest \texttt{reconcile-pending} (every 15 min) & Reconcile tests \\
FR-ADM-01 Admin session & Separate admin JWT secret, DB re-check of admin flag & Auth tests \\
FR-ADM-10 Audit log & \texttt{audit\_log} table, \texttt{lib/audit.ts} & API tests \\
FR-SEO-01/03 Sitemap, feed & \texttt{sitemap.ts}, \texttt{feed.xml} & Live: 149 URLs, 95 feed items \\
NFR-SEC-04 Rate limits & \texttt{lib/ratelimit.ts} on OTP and login & Auth tests \\
\bottomrule
\end{tabularx}
\caption{Traceability of high-priority requirements}
\label{tab:trace}
\end{table}

\section{System Architecture}

\subsection{High-Level Design}

The system is a three-tier web application split into two independently
deployed applications that share typed request/response contracts
(Figure~\ref{fig:arch}):

\begin{itemize}
  \item \textbf{Storefront and admin panel} (\texttt{code/svcar-india}):
        a Next.js application on Vercel. Catalogue, category and product
        pages are rendered on the server and cached (Incremental Static
        Regeneration); an admin change purges the affected pages through a
        webhook.
  \item \textbf{REST API} (\texttt{code/svcar-api}): a Hono service in a
        Docker container on Fly.io (Mumbai region). It owns all business
        rules, validation, pricing and database access.
  \item \textbf{Data}: PostgreSQL and object storage on Supabase; Upstash
        Redis for caching and rate limits.
  \item \textbf{External services}: PayU (payments), Resend (email), MSG91
        (order SMS), India Post (pincode lookup), Inngest (background jobs),
        Sentry (errors), Google Search and Merchant Center.
\end{itemize}

\begin{figure}[H]
\centering
\resizebox{\linewidth}{!}{%
\begin{tikzpicture}[node distance=9mm and 6mm]
  \node[ext] (cust) {Customer browser\\(storefront)};
  \node[ext, right=40mm of cust] (adm) {Owner / staff browser\\(admin panel)};
  \node[box, below=12mm of $(cust)!0.5!(adm)$, minimum width=72mm] (web) {\textbf{Vercel} -- Next.js 16 / React 19\\server-rendered, cached pages; session proxy};
  \node[box, below=10mm of web, minimum width=72mm] (api) {\textbf{Fly.io} -- Hono REST API (Node.js 20)\\JWT auth, Zod validation, pricing, orders};
  \node[ext, below left=12mm and -12mm of api] (db) {Supabase\\PostgreSQL + Storage};
  \node[ext, below=12mm of api] (payu) {PayU\\hosted checkout};
  \node[ext, below right=12mm and -12mm of api] (jobs) {Inngest, Resend,\\MSG91, Redis};
  \node[ext, left=8mm of api] (pin) {India Post\\pincode API};
  \node[ext, right=6mm of adm] (g) {Google Search /\\Merchant Center};
  \draw[arr] (cust) -- node[lbl]{HTTPS} (web);
  \draw[arr] (adm) -- node[lbl]{HTTPS} (web);
  \draw[arr] (web) -- node[lbl]{JSON / REST} (api);
  \draw[arr] (api) -- node[lbl]{Drizzle ORM} (db);
  \draw[arr,<->] (api) -- node[lbl]{hash + webhook} (payu);
  \draw[arr] (api) -- node[lbl]{events} (jobs);
  \draw[arr] (api) -- (pin);
  \draw[arr] (g) |- node[lbl,pos=0.75]{crawl, feed.xml} (web.east);
\end{tikzpicture}}
\caption{System architecture and external services}
\label{fig:arch}
\end{figure}

The browser is outside the trust boundary: totals, ownership claims and
payment status sent by the browser are never trusted. Customer session
cookies are issued through a same-origin proxy in Next.js, so they stay
first-party in browsers that block third-party cookies.

\subsection{Technical Stack}

\begin{itemize}
  \item \textbf{Framework/Language:} TypeScript throughout; Next.js 16 and
        React 19 (frontend); Hono 4 on Node.js 20 (backend).
  \item \textbf{Database:} PostgreSQL (Supabase) accessed through Drizzle ORM;
        28 tables created by 22 versioned SQL migrations.
  \item \textbf{Tools:} Git and GitHub, GitHub Actions (CI), Vitest, Mock
        Service Worker (MSW), Docker, Vercel, Fly.io, Sentry.
  \item \textbf{Libraries:} Tailwind CSS 4, Zustand (client state), SWR
        (data fetching), Zod (schema validation), bcrypt (hashing), Upstash
        Redis client, Inngest SDK.
\end{itemize}

\section{Implementation Details}

\subsection{Module 1: Core Functionality}

\paragraph{Catalogue, variants and vehicle fitment.}
Products, images and categories form the catalogue; a product can belong to
several categories through the \texttt{product\_categories} junction table.
Sellable variants are modelled as SKUs built from \emph{dimensions} (for
example finish or size) and their values, so price and stock belong to the
SKU, not to the product. Fitment links a product to vehicle models through
\texttt{product\_fitments}. The ``My Garage'' selector stores the shopper's
car in the browser and filters search and category results to parts that fit
it; product pages list every vehicle a product fits.

\paragraph{Customer and admin authentication.}
Customers sign in with a six-digit one-time password (OTP) sent by email.
The code is generated with a cryptographic random source, stored only as a
bcrypt hash, expires after five minutes and allows five attempts; the
attempt counter is incremented atomically in SQL so concurrent guesses
cannot exceed the limit. Sending is rate-limited per email and per IP
address (15 per hour). A successful login issues a short-lived access token
and an httpOnly refresh cookie whose identifier can be revoked. Customers may
later set a password. Admin tokens are signed with a separate secret, and
every admin request re-checks the admin flag in the database.

\paragraph{Cart, checkout and pricing.}
Logged-in customers have a persistent cart; guests keep the cart in the
browser and can check out without an account (their orders are linked to
the account if they later sign in with the same email). At checkout the API
reloads the SKUs from the database, rejects inactive products, and computes:

\begin{align*}
  \text{subtotal } S &= \textstyle\sum_i q_i \cdot p_i \quad(\text{current SKU prices, in paise})\\
  \text{shipping } F &= \begin{cases}0 & S \ge T_{\text{free}}\\ F_{\text{flat}} & \text{otherwise}\end{cases}\\
  \text{total } T &= (S - D) + F, \qquad \text{GST shown} = (S-D) - \frac{S-D}{1+g}
\end{align*}

where $D$ is a validated promo-code discount and $g$, $T_{\text{free}}$ and
$F_{\text{flat}}$ are admin settings (defaults: 28\% GST, free shipping from
\rupee5,000, otherwise a \rupee200 flat fee). Free shipping is decided on the
subtotal before the discount, which is the figure the shopper sees against the
threshold; all arithmetic is done in integer paise to avoid rounding errors. Prices are GST-inclusive, so
GST is displayed, not added again. For COD the customer pays a 15\% advance
online and the balance on delivery.

\paragraph{Payment processing.}
The order is created \emph{before} payment with a unique order number. The
browser is redirected to PayU with a SHA-512 signed request. The return
callback and the server-to-server webhook both verify PayU's response hash
and check the paid amount against the order before marking it paid; an
invalid callback cannot change the order. A background job runs every 15
minutes, asks PayU about orders still pending, and settles any that were paid
but whose callback was lost. A second job sends the confirmation email and SMS
when an order is paid.

\subsection{Module 2: Secondary Features}

\paragraph{Administration.}
The admin panel covers products (with a variant matrix, fitment editor and
image upload), categories (drag-and-drop ordering, archive and restore),
orders (filters, status timeline, notes, COD collection), a home-page
builder, store settings, promo codes, WhatsApp and phone leads, newsletter
subscribers, live activity, first-party analytics and the audit log.
Admins can preview a draft product through a private, non-indexed link.

\paragraph{Search engine optimisation.}
The storefront publishes an XML sitemap, a robots file, canonical URLs,
Product and Breadcrumb structured data (JSON-LD), generated Open Graph images
and a product feed for Google Merchant Center. Empty categories and admin or
preview pages are marked \texttt{noindex} so that search engines skip them.

\paragraph{White-label configuration.}
A single \texttt{brand.config.ts} file holds the store identity, contact
details, tax and shipping rules; \texttt{npm run brand:apply} writes the
derived configuration into both apps, and \texttt{brand:check} fails CI if
the generated files drift.

\paragraph{Development without a backend.}
MSW fixtures let the whole storefront and admin panel run in the browser
without the API, which the team used for UI work and for tests. Mocks are
disabled in production builds.

\subsection{Use Case Diagrams}

Figure~\ref{fig:uc} shows the actors and their goals. A registered Customer
can do everything a Guest can, plus view order history and saved addresses.

\begin{figure}[H]
\centering
\begin{tikzpicture}
  \draw[rounded corners=3pt] (-3.7,-6.4) rectangle (3.7,3.4);
  \node[font=\small\bfseries] at (0,3.05) {SV Car Customs};
  \node[uc] (u1) at (0,2.4) {Browse / search catalogue};
  \node[uc] (u2) at (0,1.6) {Check vehicle fitment};
  \node[uc] (u3) at (0,0.8) {Manage cart};
  \node[uc] (u4) at (0,0.0) {Checkout and pay};
  \node[uc] (u5) at (0,-0.8) {Sign in with email OTP};
  \node[uc] (u6) at (0,-1.6) {Track order};
  \node[uc] (u7) at (0,-2.6) {Manage products and categories};
  \node[uc] (u8) at (0,-3.4) {Process orders / collect COD};
  \node[uc] (u9) at (0,-4.2) {Edit settings and home page};
  \node[uc] (u10) at (0,-5.0) {Manage promo codes};
  \node[uc] (u11) at (0,-5.8) {View analytics and audit log};
  \actor{-5.6,0.9}{cust}{Guest /\\Customer}
  \actor{5.6,-4.2}{admin}{Admin}
  \actor{5.6,0.3}{payu}{PayU\\(external)}
  \foreach \u in {u1,u2,u3,u4,u5,u6} \draw (-5.3,1.2) -- (\u.west);
  \foreach \u in {u7,u8,u9,u10,u11} \draw (5.3,-3.9) -- (\u.east);
  \draw (5.3,0.6) -- (u4.east);
\end{tikzpicture}
\caption{Use case diagram}
\label{fig:uc}
\end{figure}

\begin{table}[H]
\centering\small
\begin{tabularx}{\linewidth}{l X X}
\toprule
\textbf{Use case} & \textbf{Precondition} & \textbf{Outcome} \\
\midrule
Checkout and pay & Non-empty cart; valid address and pincode & Order exists with a unique number; it becomes \emph{paid} only after a verified PayU response. \\
Sign in with OTP & Valid email; within rate limits & Access token and refresh cookie issued. \\
Process orders & Signed-in admin & Status advances along the timeline; change recorded in the audit log. \\
Manage products & Signed-in admin & Catalogue, variants or fitment updated; cached pages refreshed. \\
\bottomrule
\end{tabularx}
\caption{Use-case preconditions and outcomes}
\label{tab:uc}
\end{table}

\subsection{UML Class Diagrams}

Figure~\ref{fig:class} shows the main domain classes. Fitment is an
association class between Product and VehicleModel; price and stock live on
SKU. An Order keeps a snapshot of its lines and shipping address, so later
catalogue or address edits do not change past orders.

\begin{figure}[H]
\centering
\begin{tikzpicture}[font=\scriptsize,
  cls/.style={draw,rectangle split,rectangle split parts=3,fill=black!3,align=left,text width=31mm,inner sep=3pt}]
  \node[cls] (cat) at (-5,3) {\textbf{Category}\nodepart{second}- id, name, slug\\- status\nodepart{third}};
  \node[cls] (prod) at (0,3) {\textbf{Product}\nodepart{second}- id, title, slug\\- status (enum)\nodepart{third}+ fits(model): bool};
  \node[cls] (vm) at (5,3) {\textbf{VehicleModel}\nodepart{second}- id, make, model\nodepart{third}};
  \node[cls] (fit) at (5,0.6) {\textbf{Fitment}\nodepart{second}- productId, modelId\nodepart{third}};
  \node[cls] (sku) at (0,0.4) {\textbf{SKU}\nodepart{second}- id, price, stock\\- dimension values\nodepart{third}+ inStock(): bool};
  \node[cls] (img) at (-5,0.6) {\textbf{ProductImage}\nodepart{second}- url, position\nodepart{third}};
  \node[cls] (user) at (-5,-2.6) {\textbf{User}\nodepart{second}- id, email, phone\\- isAdmin\nodepart{third}+ orders()};
  \node[cls] (ord) at (0,-2.6) {\textbf{Order}\nodepart{second}- orderNumber, status\\- subtotal, discount\\- shippingFee, gst, total\\- paymentMethod\\- advanceAmount\nodepart{third}+ balanceDue()};
  \node[cls] (oi) at (5,-2.6) {\textbf{OrderItem}\nodepart{second}- skuId, title, qty\\- unitPrice (snapshot)\nodepart{third}};
  \node[cls] (ev) at (0,-6.0) {\textbf{OrderEvent}\nodepart{second}- type, note, at\nodepart{third}};
  \draw (cat) -- node[above]{0..* \quad\quad 0..*} (prod);
  \draw (prod) -- node[above]{0..* \quad\quad 0..*} (vm);
  \draw[dashed] ($(prod.east)!0.5!(vm.west)$) -- (fit);
  \draw[-{Diamond[fill=black]}] (sku) -- node[right]{0..* \ 1} (prod);
  \draw[-{Diamond[fill=black]}] (img) -- node[above]{1..* \ 1} (prod.south west);
  \draw (user) -- node[above]{1 \quad\quad 0..*} (ord);
  \draw[-{Diamond[fill=black]}] (oi) -- node[above]{1..* \quad\quad 1} (ord);
  \draw[-{Diamond[fill=black]}] (ev) -- node[right]{0..* \ 1} (ord);
\end{tikzpicture}
\caption{Domain class diagram (selected attributes)}
\label{fig:class}
\end{figure}

\subsection{UML Sequence Diagram}

Figure~\ref{fig:seq} traces an online payment. The API is the only authority
for prices and payment state.

\begin{figure}[H]
\centering
\begin{tikzpicture}[font=\scriptsize]
  \foreach \x/\n/\t in {0/c/Customer UI,3/a/Hono API,6/d/PostgreSQL,9/p/PayU,12/j/Inngest} {
    \node[box,minimum width=20mm] (\n) at (\x,0) {\t};
    \draw[dashed,black!50] (\x,-0.45) -- (\x,-9.6);
  }
  \draw[arr] (0,-1) -- node[lbl,above=1pt]{1 POST /checkout} (3,-1);
  \draw[arr] (3,-1.7) -- node[lbl,above=1pt]{2 load cart, SKUs, settings} (6,-1.7);
  \draw[arr,dashed] (6,-2.4) -- node[lbl,above=1pt]{3 trusted data} (3,-2.4);
  \draw[arr] (3,-2.9) -- ++(0.5,0) -- ++(0,-0.4) -- ++(-0.5,0); \node[lbl,right] at (3.55,-3.1) {4 compute totals};
  \draw[arr] (3,-3.9) -- node[lbl,above=1pt]{5 create order + items} (6,-3.9);
  \draw[arr,dashed] (3,-4.6) -- node[lbl,above=1pt]{6 signed PayU form} (0,-4.6);
  \draw[arr] (0,-5.3) -- node[lbl,above=1pt,pos=0.35]{7 pay (full or COD advance)} (9,-5.3);
  \draw[arr] (9,-6.0) -- node[lbl,above=1pt,pos=0.5]{8 callback / webhook} (3,-6.0);
  \draw[arr] (3,-6.5) -- ++(0.5,0) -- ++(0,-0.4) -- ++(-0.5,0); \node[lbl,right] at (3.55,-6.7) {9 verify hash + amount};
  \draw (2.2,-7.2) rectangle (12.6,-9.4); \node[anchor=north west,font=\scriptsize\bfseries,draw,fill=white] at (2.2,-7.2) {opt [verified]};
  \draw[arr] (3,-8.1) -- node[lbl,above=1pt]{10 mark paid, add event} (6,-8.1);
  \draw[arr] (3,-8.9) -- node[lbl,above=1pt,pos=0.7]{11 order.paid $\rightarrow$ email + SMS} (12,-8.9);
\end{tikzpicture}
\caption{Checkout and verified-payment sequence}
\label{fig:seq}
\end{figure}

\subsection{UML Activity Diagram}

\begin{figure}[H]
\centering
\begin{tikzpicture}[font=\scriptsize,node distance=5mm]
  \draw (-2,1) rectangle (2,-9.4); \draw (2,1) rectangle (7.6,-9.4); \draw (7.6,1) rectangle (10.6,-9.4);
  \node[font=\footnotesize\bfseries] at (0,0.7) {Customer}; \node[font=\footnotesize\bfseries] at (4.8,0.7) {Application}; \node[font=\footnotesize\bfseries] at (9.1,0.7) {PayU};
  \fill (0,0.2) circle (1mm);
  \node[box] (a1) at (0,-0.6) {Select car and SKU};
  \node[box] (a2) at (0,-1.8) {Enter address\\(sign in or guest)};
  \node[box] (a3) at (4.8,-1.8) {Validate cart, pincode};
  \node[box] (a4) at (4.8,-3.0) {Compute subtotal, promo,\\shipping, GST, total};
  \node[box] (a5) at (4.8,-4.3) {Create order (pending)};
  \node[box] (a6) at (9.1,-4.3) {Pay full /\\15\% advance};
  \node[box] (a7) at (4.8,-5.7) {Verify hash and amount};
  \node[draw,diamond,aspect=1.6,inner sep=1pt] (d) at (4.8,-6.9) {valid?};
  \node[box] (n) at (3.2,-8.3) {Keep pending,\\allow retry};
  \node[box] (y) at (6.4,-8.3) {Mark paid,\\notify};
  \draw[arr] (0,0.1) -- (a1); \draw[arr] (a1) -- (a2); \draw[arr] (a2) -- (a3); \draw[arr] (a3) -- (a4);
  \draw[arr] (a4) -- (a5); \draw[arr] (a5) -- (a6); \draw[arr] (a6) |- (a7); \draw[arr] (a7) -- (d);
  \draw[arr] (d) -| node[lbl,pos=0.3]{no} (n); \draw[arr] (d) -| node[lbl,pos=0.3]{yes} (y);
  \draw (n.south) ++(0,-0.25) circle (1.2mm); \fill (n.south) ++(0,-0.25) circle (0.7mm); \draw (n.south) -- ++(0,-0.13);
  \draw (y.south) ++(0,-0.25) circle (1.2mm); \fill (y.south) ++(0,-0.25) circle (0.7mm); \draw (y.south) -- ++(0,-0.13);
\end{tikzpicture}
\caption{Checkout activity diagram with swimlanes}
\label{fig:act}
\end{figure}

\subsection{UML State Machine Diagram}

Figure~\ref{fig:state} uses the exact states of the \texttt{order\_status}
enumeration in the database. ``Out for delivery'' is recorded as a timeline
event, not a separate status.

\begin{figure}[H]
\centering
\begin{tikzpicture}[font=\scriptsize]
  \fill (-1.7,0) circle (1mm);
  \node[state] (p) at (0,0) {pending};
  \node[state] (pa) at (3.6,0) {paid};
  \node[state] (c) at (7.2,0) {confirmed};
  \node[state] (s) at (10.8,0) {shipped};
  \node[state] (x) at (0,-2.4) {cancelled};
  \node[state] (r) at (5.4,-2.4) {refunded};
  \node[state] (d) at (10.8,-2.4) {delivered};
  \draw[arr] (-1.65,0) -- (p);
  \draw[arr] (p) -- node[lbl,above=2pt]{verified\\payment} (pa);
  \draw[arr] (pa) -- node[lbl,above=2pt]{admin\\confirms} (c);
  \draw[arr] (c) -- node[lbl,above=2pt]{dispatch} (s);
  \draw[arr] (s) -- node[lbl,right=2pt]{received} (d);
  \draw[arr] (p) -- node[lbl,left=2pt]{fail /\\cancel} (x);
  \draw[arr] (pa) -- node[lbl,left=4pt]{refund} (r);
  \draw[arr] (c) -- node[lbl,right=4pt]{refund} (r);
  \draw[arr] ([xshift=-4mm]p.north) to[out=110,in=70,looseness=3] node[lbl,above]{payment failed, retry} ([xshift=4mm]p.north);
  \draw (d.south) ++(0,-0.35) circle (1.2mm); \fill (d.south) ++(0,-0.35) circle (0.7mm); \draw (d.south) -- ++(0,-0.23);
\end{tikzpicture}
\caption{Order state machine (from the \texttt{order\_status} enum)}
\label{fig:state}
\end{figure}

\subsection{Data Flow Diagrams}

\begin{figure}[H]
\centering
\begin{tikzpicture}[font=\scriptsize]
  \node[proc,minimum size=28mm] (sys) at (0,0) {0.0\\SV Car Customs\\platform};
  \node[ext] (c) at (-6,1.8) {Customer};
  \node[ext] (a) at (6,1.8) {Admin};
  \node[ext] (p) at (-6,-1.8) {PayU};
  \node[ext] (n) at (6,-1.8) {Email / SMS\\providers};
  \node[ext] (g) at (0,-3.6) {Google};
  \draw[arr] (c.east) to[bend left=10] node[lbl,above=2pt]{search, cart,\\address, payment} (sys.north west);
  \draw[arr] (sys.west) to[bend left=10] node[lbl,below=2pt]{pages,\\order status} (c.south);
  \draw[arr,<->] (p.east) -- node[lbl,below=2pt]{payment request / result} (sys.south west);
  \draw[arr] (a.west) -- node[lbl,above=2pt]{products, settings,\\order status} (sys.north east);
  \draw[arr] (sys.south east) -- node[lbl,below=2pt]{notifications} (n.west);
  \draw[arr] (sys.south) -- node[lbl,right=2pt]{sitemap, feed} (g.north);
\end{tikzpicture}
\caption{Level-0 data flow diagram (context)}
\label{fig:dfd0}
\end{figure}

\begin{figure}[H]
\centering
\begin{tikzpicture}[font=\scriptsize]
  \node[ext] (c) at (-5,2.8) {Customer};
  \node[ext] (p) at (0,2.8) {PayU};
  \node[store] (d3) at (2.6,2.8) {D3 Settings};
  \node[ext] (a) at (5.2,2.8) {Admin};
  \node[proc] (p1) at (-5,0) {1.0\\Browse\\and cart};
  \node[proc] (p2) at (0,0) {2.0\\Checkout\\and payment};
  \node[proc] (p3) at (5.2,0) {3.0\\Order\\management};
  \node[store] (d1) at (-5,-3) {D1 Catalogue};
  \node[store] (d2) at (2.6,-3) {D2 Orders};
  \draw[arr] (c) -- node[lbl]{search,\\add to cart} (p1);
  \draw[arr] (d1) -- node[lbl]{products,\\fitment} (p1);
  \draw[arr] (p1) -- node[lbl,above=2pt]{cart items} (p2);
  \draw[arr,<->] (p) -- node[lbl]{payment} (p2);
  \draw[arr] (d3.south west) -- node[lbl,pos=0.45]{GST,\\shipping} (p2.north east);
  \draw[arr] (p3.north west) -- node[lbl,pos=0.55]{settings} (d3.south east);
  \draw[arr] (p2.south east) -- node[lbl]{new order} (d2.north west);
  \draw[arr,<->] (d2.north east) -- node[lbl]{order updates} (p3.south west);
  \draw[arr,<->] (a) -- node[lbl]{status,\\settings} (p3);
\end{tikzpicture}
\caption{Level-1 data flow diagram}
\label{fig:dfd1}
\end{figure}

\subsection{Data Model}

The database has 28 tables (Table~\ref{tab:data}). Customer and Admin are
roles of one \texttt{users} table with an \texttt{is\_admin} flag.

\begin{table}[H]
\centering\footnotesize
\begin{tabularx}{\linewidth}{l X}
\toprule
\textbf{Area} & \textbf{Tables} \\
\midrule
Catalogue & categories, products, product\_images, product\_categories \\
Variants & dimensions, dimension\_values, skus \\
Fitment & vehicle\_makes, vehicle\_models, product\_fitments \\
Customers & users, otp\_codes, addresses \\
Commerce & carts, cart\_items, orders, order\_items, order\_events, order\_backups, promo\_codes \\
Operations & settings, home\_config, audit\_log, admin\_push\_subscriptions \\
Growth & subscribers, analytics\_events, whatsapp\_leads, phone\_leads \\
\bottomrule
\end{tabularx}
\caption{Data model by area (28 tables)}
\label{tab:data}
\end{table}

% =====================================================================
\chapter{Results and Evaluation}\label{ch:results}
% =====================================================================

\section{Testing Methodology}\label{sec:testing}

\begin{itemize}
  \item \textbf{Unit testing:} Vitest suites cover pricing, promo rules,
        payment hashes, payment amounts, OTP limits, cart state, checkout
        helpers, category trees, SEO metadata and UI components.
  \item \textbf{Integration testing:} API route tests run full HTTP requests
        through the Hono app with test doubles for the database and external
        services; storefront tests run against MSW handlers that mimic the
        API contracts. A separate end-to-end script places a real checkout
        against a running API.
  \item \textbf{Security testing:} a test catalogue mapped to OWASP risks
        (ownership checks, JWT handling, payment integrity, injection,
        uploads, business-logic abuse) guided dedicated tests, for example
        that a callback with a wrong hash or amount cannot mark an order paid.
  \item \textbf{Continuous integration:} every pull request runs API
        type-check, tests and build, storefront type-check, tests, build and
        lint, and the brand-drift check. Separate workflows block direct
        pushes to \texttt{main} and flag changes to sensitive paths.
  \item \textbf{Performance testing:} response times were measured on the
        live site and Lighthouse (mobile profile) was run on live pages.
  \item \textbf{User testing:} no formal usability study was run; real
        usage was observed through analytics after launch (Section~\ref{sec:analytics}).
\end{itemize}

\begin{table}[H]
\centering\small
\begin{tabular}{l r r r}
\toprule
\textbf{Suite} & \textbf{Test files} & \textbf{Tests} & \textbf{Result} \\
\midrule
API (\texttt{svcar-api}) & 56 & 399 & 398 passed, 1 skipped \\
Storefront (\texttt{svcar-india}) & 57 & 486 & 486 passed \\
Type-check (both apps) & -- & -- & no errors \\
\midrule
\textbf{Total} & \textbf{113} & \textbf{885} & \textbf{884 passed, 0 failed} \\
\bottomrule
\end{tabular}
\caption{Automated test results (run on 26 November 2026)}
\label{tab:tests}
\end{table}

The skipped API test needs a live database and is run only in that
environment.

\section{Performance Metrics}

\begin{table}[H]
\centering\small
\begin{tabularx}{\linewidth}{X l l}
\toprule
\textbf{Metric} & \textbf{Target (SRS)} & \textbf{Achieved} \\
\midrule
Cached page, time to first byte (home, shop, category; median of 3) & $<$ 300\,ms & 375--437\,ms \\
Sitemap / feed response & -- & 334\,ms / 545\,ms \\
API response, public endpoint (median of 5) & $<$ 1\,s & 1.26\,s \\
LCP, Lighthouse mobile (home / category) & $<$ 2.5\,s & 10.8\,s / 9.4\,s \\
Lighthouse SEO score & -- & 100 / 100 \\
Lighthouse accessibility score & -- & 91 / 94 \\
Lighthouse performance score & -- & 62 / 63 \\
Cumulative Layout Shift & -- & 0 / 0 \\
Automated tests passing & 100\% & 100\% (884 of 884 run) \\
Indexed URLs in sitemap / items in feed & -- & 149 / 95 \\
\bottomrule
\end{tabularx}
\caption{Measured results against targets (live site, 26 November 2026)}
\label{tab:perf}
\end{table}

Response times were measured from a client in India over a home connection,
so they include network latency. Lighthouse simulates a slow 4G mobile
connection and a throttled CPU; its LCP is a pessimistic lab figure, not the
75th-percentile field value the SRS target refers to, which needs Vercel
Speed Insights data that is not yet enabled.

\section{Post-launch Usage}\label{sec:analytics}

\begin{table}[H]
\centering\small
\begin{tabularx}{\linewidth}{X l}
\toprule
\textbf{Google Search Console, 15 Oct -- 11 Nov 2026} & \textbf{Value} \\
\midrule
Clicks from Google search & 81 \\
Impressions & 2,477 \\
Click-through rate & 3.3\% \\
Average position & 14.8 \\
Top query by clicks & ``svcarcustoms'' (13 clicks, position 1.9) \\
\bottomrule
\end{tabularx}
\caption{Organic search performance after launch}
\label{tab:gsc}
\end{table}

The store's own analytics for the 28 days before 9 November 2026 recorded
about 850 visitors (Instagram and direct visits were the largest sources),
71 checkout views and 1 confirmed order. These first-party numbers do not yet
filter out bots or the team's own visits.

\section{Analysis}

\paragraph{Did we meet our objectives?}
\begin{itemize}
  \item \textbf{Objective 1 (fitment-aware shopping): met.} Fitment, ``My
        Garage'' filtering, guest checkout and a single checkout page (cart $\rightarrow$
        checkout $\rightarrow$ PayU, within the three-screen limit) are live.
  \item \textbf{Objective 2 (owner control): met.} Products, prices, tax,
        shipping, home page and promo codes are all edited in the admin panel;
        every change is written to the audit log.
  \item \textbf{Objective 3 (payment integrity): met.} Hash and amount
        verification and the 15-minute reconciliation job are implemented and
        covered by tests.
  \item \textbf{Objective 4 (search visibility): met for the technical part.}
        Sitemap, structured data and feed are live and Lighthouse SEO scores
        100; organic traffic is still mostly brand searches.
  \item \textbf{Objective 5 (quality): partly met.} CI and the test suite are
        in place and green, but measured page speed (Table~\ref{tab:perf})
        does not yet meet the 300\,ms and 2.5\,s targets.
\end{itemize}

\paragraph{Where did we exceed expectations?}
Scope grew beyond the SRS: guest checkout, promo codes, first-party analytics,
admin push notifications, phone leads and the white-label brand
configuration were not in version~1.0 of the SRS.

\paragraph{What challenges did we encounter, and how were they resolved?}
\begin{itemize}
  \item \emph{Third-party cookies:} refresh cookies set by the API domain were
        blocked by some browsers, logging customers out. Session routes were
        moved behind a same-origin Next.js proxy so the cookie is first-party.
  \item \emph{Image optimisation quota:} Vercel's free image-optimisation
        quota ran out; a custom image loader now uses Supabase's image
        rendering endpoint.
  \item \emph{Database connectivity outage:} the API could not reach the
        database over a direct IPv6 connection, which took the store offline.
        Switching to Supabase's session pooler restored service; cached
        storefront pages are kept (not replaced by error pages) while the API is down.
  \item \emph{Low checkout conversion:} analytics showed many checkout views
        but few orders; mobile layout fixes were shipped and conversion is still
        being monitored.
\end{itemize}

% =====================================================================
\chapter{Conclusion}
% =====================================================================

\section{Summary of Work}

The team built and deployed a complete vehicle-aware e-commerce platform for
a real business: a Next.js storefront and admin panel, a Hono REST API, a
28-table PostgreSQL schema and integrations with PayU, Inngest, Resend,
MSG91, India Post and Google. The work was delivered in weekly increments,
documented through an SRS, UML views, data-flow diagrams and a data model,
and verified by 884 passing automated tests in CI. Four of the five objectives
were met in full; the performance objective was met for testing and CI but not
yet for page speed.

\section{Team Contributions}

The work was divided into four equal areas, one per member.

\begin{table}[H]
\centering\small
\begin{tabularx}{\linewidth}{l l X}
\toprule
\textbf{Member} & \textbf{Area} & \textbf{Main work} \\
\midrule
Vishu Verma & Payments and deployment & PayU integration, COD advance, payment reconciliation job, checkout API, Vercel and Fly.io deployment \\
Raghu Verma & Backend and database & Database schema and migrations, REST API routes, email-OTP and admin authentication, rate limits, audit log \\
Pratham Sharma & Frontend & Storefront pages, admin panel UI, My Garage selector, image loading, mobile layout \\
Tarush Mohindru & Quality and SEO & Test suites, CI workflows, SEO (sitemap, structured data, feed), SRS, UML diagrams and reports \\
\bottomrule
\end{tabularx}
\caption{Division of work}
\label{tab:team}
\end{table}

\section{Key Findings}

\begin{itemize}
  \item \textbf{Finding 1:} Modelling fitment as a product--vehicle
        association and price/stock on the SKU kept variant selection and
        compatibility independent, which made the ``My Garage'' filter a
        query rather than a special case.
  \item \textbf{Finding 2:} Treating the browser as untrusted (server-side
        pricing, hash and amount verification, reconciliation) was essential;
        payment callbacks do get lost in practice, and the reconciliation job
        is what guarantees no paid order is lost.
  \item \textbf{Finding 3:} The unexpected result was that technical SEO was
        not enough on its own: Lighthouse SEO scored 100, but almost all Google
        clicks still came from brand searches. Ranking for product queries
        needs product-specific content and time.
\end{itemize}

\section{Future Work and Recommendations}

\begin{enumerate}
  \item Verified-purchase customer reviews and ratings (designed and
        implemented on a separate branch; awaiting database migration).
  \item Automatic courier booking and live shipment tracking.
  \item Page-speed work towards the 2.5\,s LCP target: smaller hero images,
        fewer client-side scripts on the home page, and field measurement
        with Vercel Speed Insights.
  \item Moving the API and database to the same region to remove
        cross-region latency from every API call.
  \item A small usability study (5--8 users) of the vehicle-selection and
        checkout flow, and a native Android app.
\end{enumerate}

\section{Final Remarks}

Building a system that real customers pay through taught us that software
engineering practice matters most at the boundaries: requirements that can be
traced to tests, a CI pipeline that blocks unsafe changes, and designs that
assume external services will fail. The platform now runs the business's
online sales, and the remaining work is measured and prioritised rather than
guessed.

\begin{thebibliography}{PayU25}
\bibitem[Next25]{next} Vercel. \emph{Next.js Documentation (App Router)}. \url{https://nextjs.org/docs}, 2025.
\bibitem[Hono25]{hono} Hono. \emph{Hono -- Web framework built on Web Standards}. \url{https://hono.dev}, 2025.
\bibitem[Driz25]{drizzle} Drizzle Team. \emph{Drizzle ORM Documentation}. \url{https://orm.drizzle.team}, 2025.
\bibitem[PayU25]{payu} PayU India. \emph{PayU Hosted Checkout Integration}. \url{https://docs.payu.in}, 2025.
\bibitem[OWASP21]{owasp} OWASP Foundation. \emph{OWASP Top 10: 2021}. \url{https://owasp.org/Top10}, 2021.
\bibitem[IEEE830]{ieee830} IEEE. \emph{IEEE Recommended Practice for Software Requirements Specifications (Std 830-1998)}. 1998.
\bibitem[Goog25]{google} Google. \emph{Search Central: Product structured data}. \url{https://developers.google.com/search/docs}, 2025.
\bibitem[SRS26]{srs} V.~Verma, R.~Verma, P.~Sharma, T.~Mohindru. \emph{Software Requirements Specification for SV Car Customs, v1.0}. UCS503, TIET, 15 Oct 2026.
\end{thebibliography}

\end{document}
